\documentclass[11pt]{article}

\usepackage[utf8]{inputenc}
\usepackage[T1]{fontenc}
\usepackage{lmodern}
\usepackage{geometry}
\usepackage{amsmath,amssymb,amsfonts,amsthm,mathtools}
\usepackage{bm}
\usepackage{enumitem}
\usepackage{microtype}
\usepackage{hyperref}
\usepackage{titlesec}
\usepackage{booktabs}
\usepackage{array}
\usepackage{setspace}
\usepackage{mathrsfs}
\usepackage{physics}

\geometry{margin=1in}
\hypersetup{colorlinks=true, linkcolor=black, urlcolor=blue, citecolor=black}
\setlist[enumerate]{topsep=0.4em,itemsep=0.35em}
\setlist[itemize]{topsep=0.3em,itemsep=0.25em}
\titleformat{\section}{\large\bfseries}{\thesection.}{0.5em}{}
\titleformat{\subsection}{\normalsize\bfseries}{\thesubsection.}{0.5em}{}
\setstretch{1.06}

\newcommand{\Aether}{\AE{}ther}
\newcommand{\Aflow}{\AE{}ther-flow}
\newcommand{\TheoryName}{\Aflow{} Interpretation of Relativity}
\newcommand{\TheoryProgram}{\Aether{} / \Aflow{} framework}

\newcommand{\CompletedMetric}{g^{\sharp}}
\newcommand{\MatterSpace}{\Psi}

\newcommand{\Car}{\mathrm{car}}
\newcommand{\Cur}{\mathrm{cur}}
\newcommand{\Hom}{\mathrm{hom}}

\newcommand{\ForSpace}[1]{\mathcal I_{#1}^{\mathrm{for}}}
\newcommand{\PrimShell}{\mathcal S_{\mathrm{prim}}}
\newcommand{\MetricOut}{\mathscr G_{\mathrm{out}}}
\newcommand{\MatterOut}{\mathscr T_{\mathrm{out}}}

\newcommand{\SubSectionPackage}{\mathfrak V^{\mathrm{subsec}}}
\newcommand{\MicroSectionPackage}{\mathfrak W^{\mathrm{microsec}}}
\newcommand{\MicroSectionSpace}[1]{\mathcal N_{#1}^{\mathrm{microsec}}}
\newcommand{\MicroSectionBody}[1]{\mathcal B_{#1}^{\mathrm{microsec}}}
\newcommand{\MicroSectionFunctional}[1]{\mathscr E_{#1}^{\mathrm{microsec}}}
\newcommand{\MicroSectionShell}[1]{\mathcal Z_{#1}^{\mathrm{microsec}}}
\newcommand{\SubPreSectionCore}{\mathfrak U^{\mathrm{sec}}}

\newcommand{\NanoSectionPackage}{\mathfrak X^{\mathrm{nanosec}}}
\newcommand{\NanoSectionSpace}[1]{\mathcal P_{#1}^{\mathrm{nanosec}}}
\newcommand{\NanoSectionBody}[1]{\mathcal B_{#1}^{\mathrm{nanosec}}}
\newcommand{\NanoSectionProj}[1]{\Pi_{#1}^{\mathrm{nanosec}}}
\newcommand{\NanoSectionFiber}[2]{\mathcal F_{#1}^{\mathrm{nanosec}}(#2)}
\newcommand{\NanoSectionMicroMin}[1]{\mathfrak n_{#1}^{\mathrm{nanosec}}}
\newcommand{\NanoSectionFunctional}[1]{\mathscr N_{#1}^{\mathrm{nanosec}}}
\newcommand{\NanoSectionShell}[1]{\mathcal Z_{#1}^{\mathrm{nanosec}}}

\newcommand{\PicoSectionPackage}{\mathfrak Y^{\mathrm{picosec}}}
\newcommand{\PicoSectionSpace}[1]{\mathcal Q_{#1}^{\mathrm{picosec}}}
\newcommand{\PicoSectionBody}[1]{\mathcal B_{#1}^{\mathrm{picosec}}}
\newcommand{\PicoSectionProj}[1]{\Pi_{#1}^{\mathrm{picosec}}}
\newcommand{\PicoSectionFiber}[2]{\mathcal F_{#1}^{\mathrm{picosec}}(#2)}
\newcommand{\PicoSectionNanoMin}[1]{\mathfrak p_{#1}^{\mathrm{picosec}}}
\newcommand{\PicoSectionFunctional}[1]{\mathscr P_{#1}^{\mathrm{picosec}}}
\newcommand{\PicoSectionShell}[1]{\mathcal Z_{#1}^{\mathrm{picosec}}}

\newcommand{\PicoImageCore}{\mathfrak Y^{\mathrm{pic,img}}}
\newcommand{\PicoImgBody}[1]{\mathcal B_{#1}^{\mathrm{pic,img}}}
\newcommand{\PicoImgProj}[1]{\Pi_{#1}^{\mathrm{pic,img}}}
\newcommand{\PicoImgFunctional}[1]{\mathscr P_{#1}^{\mathrm{pic,img}}}
\newcommand{\PicoImgShell}[1]{\mathcal Z_{#1}^{\mathrm{pic,img}}}

\newcommand{\PicoMinSectionCore}{\mathfrak Y^{\mathrm{pic,sec}}}
\newcommand{\PicoMinSection}[1]{\mathfrak s_{#1}^{\mathrm{pic,sec}}}

\newtheorem{definition}{Definition}
\newtheorem{proposition}{Proposition}
\newtheorem{remark}{Remark}
\newtheorem{corollary}{Corollary}
\newtheorem{theorem}{Theorem}

\title{\textbf{The \TheoryName{}}\\[0.4em]
\large Primitive-Reservoir Observer-Reduction\\
\large Observer-Relevant Pico-Section Minimizer-Image Core\\
\large Collapse to an Observer-Relevant Pico-Section Minimizer-Section Core}
\author{}
\date{}

\begin{document}

\maketitle

\begin{abstract}
The current primitive-reservoir observer-reduction frontier already compressed the live pico-section burden from the full nano-section-generating substrate pico-section dynamics package \(\PicoSectionPackage\) to the smaller observer-relevant pico-section minimizer-image core \(\PicoImageCore\). Under the recorded safeguard against recursive depth drift, the next honest move below that frontier is therefore not another deeper-origin relay that preserves the same image core. The sharper benchmark-facing question is whether the image core itself still contains presentational redundancy once the recorded nano-section package is treated as fixed.

The present note proves that it does. Relative to the already fixed nano-section body, nano-section functional, and exact nano-section shell, the only independent pico-section datum still carried by the current image core is the canonical minimizer section over the recorded nano-section body. The minimizer-image body is exactly the image of that section, the restricted projection is its inverse, the exact-shell image is its shell restriction, and the restricted image-core functional is canonically induced from the recorded nano-section functional along the same section. Separate listing of those image-core constituents therefore does not add independent benchmark-facing freedom.

The benchmark boundary remains fixed throughout. The one operative metric is still exactly \(\CompletedMetric\), the ordinary matter route is unchanged, there is no second operative metric, the low-energy matter law is not enlarged, and no stronger promotion language is licensed. The positive result is therefore a genuine smaller theorem-level collapse strictly inside the current image core: the live burden now contracts to an observer-relevant pico-section minimizer-section core \(\PicoMinSectionCore\). What remains open is to derive or justify that section core from still more primitive substrate structure, or else prove a still sharper collapse strictly inside it.
\end{abstract}

\section{Introduction}

The active derivational program of \TheoryName{} again reaches a stage where depth alone is not enough. The current active-primary theorem already showed that the full nano-section-generating substrate pico-section dynamics package \(\PicoSectionPackage\) is not the minimal benchmark-facing datum at present scope, because the live burden there collapses to the smaller observer-relevant pico-section minimizer-image core \(\PicoImageCore\). After that step, another deeper-origin relay that merely preserves \(\PicoImageCore\) is no longer the honest default move.

That theorem also exposes the next honest question. Inside the image core itself, what is genuinely independent observer-relevant data, and what is merely presentational bookkeeping once the recorded nano-section package is already fixed? The present note takes the sharper internal-collapse route rather than a same-output deeper relay. It shows that the image-core body, the restricted projection, the shell image, and the restricted image-core functional are not independent pieces at current scope. They are all recovered from a single canonical minimizer section over the recorded nano-section body.

\paragraph{Framework and claim boundary.}
\TheoryProgram{} denotes the broader \Aether{} / \Aflow{} framework built on the claims that the \Aether{} is the underlying four-dimensional substrate of reality, the \Aflow{} is its intrinsic ordered motion, observed three-dimensional space is the local experiential slice of that deeper substrate, S-time is the experienced order of change arising from matter, light, and the \Aflow{}, and observed expansion is the three-dimensional appearance of deeper four-dimensional ordered motion. Within that framework, \TheoryName{} denotes the exact-closure theory in which the effective gravitational dynamics are adopted to be exactly Einsteinian with universal matter coupling. In the active sequence, \emph{adoption} means use of that established relativistic dynamics without claiming substrate derivation, while \emph{derivation} is reserved for a first-principles recovery from explicit substrate variables.

The present note stays strictly inside that boundary. It does not introduce a second operative metric, it does not enlarge the low-energy matter law, and it does not reopen the already-screened full pico-section package or the larger transport-image, sub-pre-core, or pre-core presentations as if they were again the live burden. Its narrower benchmark-facing role is to isolate the only remaining independent datum still carried by the current pico-section minimizer-image core.

\section{Fixed Inputs}

\begin{definition}[Recorded primitive continuation data]
\label{def:fixed}
Fix the following continuation data from the active primitive-reservoir observer-reduction line.
\begin{enumerate}
    \item For each sector label \(\sigma\in\{\Car,\Cur,\Hom\}\), a primitive forcing shell
    \[
    \ForSpace{\sigma}.
    \]
    \item The product shell
    \[
    \PrimShell
    \equiv
    \ForSpace{\Car}\times \ForSpace{\Cur}\times \ForSpace{\Hom}.
    \]
    \item The recorded metric and ordinary-matter outputs
    \[
    \MetricOut:\PrimShell\to\{\text{observer metrics}\},
    \qquad
    \MatterOut:\PrimShell\times\MatterSpace\to\{\text{observer stress tensors}\},
    \]
    satisfying
    \begin{equation}
    \MetricOut(\mathbf w)=\CompletedMetric,
    \qquad
    \MatterOut(\mathbf w,\psi)
    =
    T_{\mu\nu}^{\mathrm{matter}}[\CompletedMetric,\psi]
    \label{eq:fixedoutputs}
    \end{equation}
    for every \(\mathbf w\in\PrimShell\) and every matter configuration \(\psi\in\MatterSpace\).
\end{enumerate}
\end{definition}

\begin{definition}[Recorded nano-section downstream chain]
\label{def:downstream}
Fix \(\sigma\in\{\Car,\Cur,\Hom\}\). The active line already fixes the following downstream data and consequences:
\begin{enumerate}
    \item the recorded micro-section target data
    \[
    \MicroSectionSpace{\sigma},\qquad
    \MicroSectionBody{\sigma},\qquad
    \MicroSectionFunctional{\sigma},\qquad
    \MicroSectionShell{\sigma};
    \]
    \item the current exact-shell-transport-section-generating substrate sub-section dynamics package \(\SubSectionPackage\), the current micro-section package \(\MicroSectionPackage\), and the observer-relevant sub-pre-core exact-shell transport-section core \(\SubPreSectionCore\);
    \item the previously recorded fact that, once the recorded nano-section package is available, the current micro-section package, the current sub-section package, the observer-relevant section core, and the downstream shell-transport consequences used on the active line follow unchanged.
\end{enumerate}
\end{definition}

\begin{definition}[Recorded micro-section-generating substrate nano-section dynamics]
\label{def:nanosec}
Fix \(\sigma\in\{\Car,\Cur,\Hom\}\). The recorded current nano-section dynamics package \(\NanoSectionPackage\) for sector \(\sigma\) consists of the following data.
\begin{enumerate}
    \item A finite-dimensional Euclidean nano-section state space \(\NanoSectionSpace{\sigma}\), a compact convex nano-section body
    \[
    \NanoSectionBody{\sigma}\subset\NanoSectionSpace{\sigma}
    \]
    with nonempty interior, and an affine surjection
    \[
    \NanoSectionProj{\sigma}:\NanoSectionSpace{\sigma}\to\MicroSectionSpace{\sigma}
    \]
    satisfying
    \[
    \NanoSectionProj{\sigma}\bigl(\NanoSectionBody{\sigma}\bigr)=\MicroSectionBody{\sigma}.
    \]
    For each \(u\in\MicroSectionBody{\sigma}\), write
    \[
    \NanoSectionFiber{\sigma}{u}
    \equiv
    \bigl(\NanoSectionProj{\sigma}\bigr)^{-1}(u)\cap\NanoSectionBody{\sigma}.
    \]
    Each such fiber is required to be nonempty, compact, and convex.
    \item A nonnegative smooth nano-section functional
    \[
    \NanoSectionFunctional{\sigma}:\NanoSectionSpace{\sigma}\to[0,\infty)
    \]
    such that, for every \(u\in\MicroSectionBody{\sigma}\), the restriction of \(\NanoSectionFunctional{\sigma}\) to \(\NanoSectionFiber{\sigma}{u}\) is strictly convex and attains a unique minimizer. The resulting minimizer depends smoothly on \(u\); write it as
    \[
    \NanoSectionMicroMin{\sigma}:\MicroSectionBody{\sigma}\to\NanoSectionBody{\sigma}.
    \]
    These data satisfy
    \[
    \NanoSectionProj{\sigma}\circ\NanoSectionMicroMin{\sigma}
    =
    \mathrm{id}_{\MicroSectionBody{\sigma}},
    \]
    and the effective minimum-value map recovers the recorded micro-section functional:
    \[
    \MicroSectionFunctional{\sigma}(u)
    =
    \NanoSectionFunctional{\sigma}\bigl(\NanoSectionMicroMin{\sigma}(u)\bigr)
    =
    \min_{p\in\NanoSectionFiber{\sigma}{u}}
    \NanoSectionFunctional{\sigma}(p)
    \qquad
    \text{for every }
    u\in\MicroSectionBody{\sigma}.
    \]
    \item A closed embedded exact nano-section shell
    \[
    \NanoSectionShell{\sigma}\subset\NanoSectionBody{\sigma}
    \]
    satisfying
    \[
    \NanoSectionShell{\sigma}
    =
    \left\{
    p\in\NanoSectionBody{\sigma}:
    \begin{array}{l}
    \NanoSectionProj{\sigma}(p)\in\MicroSectionShell{\sigma},\\[0.2em]
    \NanoSectionFunctional{\sigma}(p)
    =
    \displaystyle\min_{q\in\NanoSectionFiber{\sigma}{\NanoSectionProj{\sigma}(p)}}
    \NanoSectionFunctional{\sigma}(q)
    \end{array}
    \right\},
    \]
    and
    \[
    \NanoSectionProj{\sigma}\big|_{\NanoSectionShell{\sigma}}
    :
    \NanoSectionShell{\sigma}
    \xrightarrow{\cong}
    \MicroSectionShell{\sigma}
    \]
    is a diffeomorphism.
    \item Benchmark faithfulness, meaning preservation of the benchmark outputs \eqref{eq:fixedoutputs}, no second operative metric, no enlarged low-energy matter law, and presentation as the fixed benchmark-relevant input for the current downstream chain.
\end{enumerate}
\end{definition}

\begin{definition}[Recorded pico-section dynamics]
\label{def:picosec}
Fix \(\sigma\in\{\Car,\Cur,\Hom\}\). The recorded current pico-section dynamics package \(\PicoSectionPackage\) for sector \(\sigma\) consists of the following data.
\begin{enumerate}
    \item A finite-dimensional Euclidean pico-section state space \(\PicoSectionSpace{\sigma}\), a compact convex pico-section body
    \[
    \PicoSectionBody{\sigma}\subset\PicoSectionSpace{\sigma}
    \]
    with nonempty interior, and an affine surjection
    \[
    \PicoSectionProj{\sigma}:\PicoSectionSpace{\sigma}\to\NanoSectionSpace{\sigma}
    \]
    satisfying
    \[
    \PicoSectionProj{\sigma}\bigl(\PicoSectionBody{\sigma}\bigr)=\NanoSectionBody{\sigma}.
    \]
    For each \(p\in\NanoSectionBody{\sigma}\), write
    \[
    \PicoSectionFiber{\sigma}{p}
    \equiv
    \bigl(\PicoSectionProj{\sigma}\bigr)^{-1}(p)\cap\PicoSectionBody{\sigma}.
    \]
    Each such fiber is required to be nonempty, compact, and convex.
    \item A nonnegative smooth pico-section functional
    \[
    \PicoSectionFunctional{\sigma}:\PicoSectionSpace{\sigma}\to[0,\infty)
    \]
    such that, for every \(p\in\NanoSectionBody{\sigma}\), the restriction of \(\PicoSectionFunctional{\sigma}\) to \(\PicoSectionFiber{\sigma}{p}\) is strictly convex and attains a unique minimizer. The resulting minimizer depends smoothly on \(p\); write it as
    \[
    \PicoSectionNanoMin{\sigma}:\NanoSectionBody{\sigma}\to\PicoSectionBody{\sigma}.
    \]
    These data satisfy
    \[
    \PicoSectionProj{\sigma}\circ\PicoSectionNanoMin{\sigma}
    =
    \mathrm{id}_{\NanoSectionBody{\sigma}},
    \]
    and the effective minimum value recovers the recorded nano-section functional:
    \[
    \NanoSectionFunctional{\sigma}(p)
    =
    \PicoSectionFunctional{\sigma}\bigl(\PicoSectionNanoMin{\sigma}(p)\bigr)
    =
    \min_{r\in\PicoSectionFiber{\sigma}{p}}
    \PicoSectionFunctional{\sigma}(r)
    \qquad
    \text{for every }
    p\in\NanoSectionBody{\sigma}.
    \]
    \item A closed embedded exact pico-section shell
    \[
    \PicoSectionShell{\sigma}\subset\PicoSectionBody{\sigma}
    \]
    satisfying
    \[
    \PicoSectionShell{\sigma}
    =
    \left\{
    r\in\PicoSectionBody{\sigma}:
    \begin{array}{l}
    \PicoSectionProj{\sigma}(r)\in\NanoSectionShell{\sigma},\\[0.2em]
    \PicoSectionFunctional{\sigma}(r)
    =
    \displaystyle\min_{s\in\PicoSectionFiber{\sigma}{\PicoSectionProj{\sigma}(r)}}
    \PicoSectionFunctional{\sigma}(s)
    \end{array}
    \right\},
    \]
    and
    \[
    \PicoSectionProj{\sigma}\big|_{\PicoSectionShell{\sigma}}
    :
    \PicoSectionShell{\sigma}
    \xrightarrow{\cong}
    \NanoSectionShell{\sigma}
    \]
    is a diffeomorphism.
    \item Benchmark faithfulness, meaning preservation of the benchmark outputs \eqref{eq:fixedoutputs}, no second operative metric, no enlarged low-energy matter law, and presentation as a deeper substrate-side source for the current nano-section package rather than as a replacement benchmark theory.
\end{enumerate}
\end{definition}

\begin{remark}
Definitions \ref{def:nanosec} and \ref{def:picosec} fix the current deeper package below the nano-section line. The previous active-primary theorem already showed that benchmark-facingly relevant pico-section data factor through a smaller minimizer-image core.
\end{remark}

\section{Current Observer-Relevant Image-Core Frontier}

\begin{definition}[Observer-relevant pico-section minimizer-image core]
\label{def:imagecore}
Fix \(\sigma\in\{\Car,\Cur,\Hom\}\) and a benchmark-faithful pico-section package in the sense of Definition \ref{def:picosec}. The \emph{observer-relevant pico-section minimizer-image core} \(\PicoImageCore\) for sector \(\sigma\) consists of:
\begin{enumerate}
    \item the compact minimizer-image body
    \[
    \PicoImgBody{\sigma}
    \equiv
    \PicoSectionNanoMin{\sigma}\bigl(\NanoSectionBody{\sigma}\bigr)
    \subset
    \PicoSectionBody{\sigma},
    \]
    together with the restricted projection
    \[
    \PicoImgProj{\sigma}
    \equiv
    \PicoSectionProj{\sigma}\big|_{\PicoImgBody{\sigma}}
    :
    \PicoImgBody{\sigma}\to\NanoSectionBody{\sigma};
    \]
    \item the restricted functional
    \[
    \PicoImgFunctional{\sigma}
    \equiv
    \PicoSectionFunctional{\sigma}\big|_{\PicoImgBody{\sigma}};
    \]
    \item the exact-shell image
    \[
    \PicoImgShell{\sigma}
    \equiv
    \PicoSectionNanoMin{\sigma}\bigl(\NanoSectionShell{\sigma}\bigr)
    \subset
    \PicoSectionShell{\sigma};
    \]
    \item benchmark faithfulness in the sense of Definition \ref{def:picosec}(4).
\end{enumerate}
\end{definition}

\begin{remark}
Definition \ref{def:imagecore} is the current benchmark-facing frontier before the present theorem is applied. The question now is whether its body, projection, functional, and shell constituents are all independently live, or whether some of them are already forced once the recorded nano-section package is fixed.
\end{remark}

\section{Observer-Relevant Pico-Section Minimizer-Section Core}

\begin{definition}[Observer-relevant pico-section minimizer-section core]
\label{def:sectioncore}
Fix \(\sigma\in\{\Car,\Cur,\Hom\}\) and a benchmark-faithful pico-section package in the sense of Definition \ref{def:picosec}. The \emph{observer-relevant pico-section minimizer-section core} \(\PicoMinSectionCore\) for sector \(\sigma\) consists of:
\begin{enumerate}
    \item the canonical minimizer section
    \[
    \PicoMinSection{\sigma}
    \equiv
    \PicoSectionNanoMin{\sigma}
    :
    \NanoSectionBody{\sigma}\to\PicoSectionBody{\sigma};
    \]
    \item benchmark faithfulness in the sense of Definition \ref{def:picosec}(4).
\end{enumerate}
The image body \(\PicoMinSection{\sigma}\bigl(\NanoSectionBody{\sigma}\bigr)\), the shell restriction
\[
\PicoMinSection{\sigma}\big|_{\NanoSectionShell{\sigma}}
:
\NanoSectionShell{\sigma}\to\PicoSectionShell{\sigma},
\]
the inverse projection on that image, and the image-core functional induced from \(\NanoSectionFunctional{\sigma}\) are understood as derived data rather than separately listed core data.
\end{definition}

\begin{remark}
Definition \ref{def:sectioncore} is strictly smaller than Definition \ref{def:imagecore} at current scope. Once the recorded nano-section body, nano-section functional, and exact nano-section shell are fixed, the image body, restricted projection, shell image, and restricted image-core functional are all recovered from the single canonical minimizer section.
\end{remark}

\section{Collapse of the Image Core}

\begin{proposition}[The canonical minimizer section reconstructs the image body and projection]
\label{prop:reconstructbody}
For every benchmark-faithful pico-section package,
\[
\PicoImgBody{\sigma}
=
\PicoMinSection{\sigma}\bigl(\NanoSectionBody{\sigma}\bigr),
\]
and the restricted projection
\[
\PicoImgProj{\sigma}:\PicoImgBody{\sigma}\to\NanoSectionBody{\sigma}
\]
is a homeomorphism with inverse \(\PicoMinSection{\sigma}\).
\end{proposition}

\begin{proof}
The displayed identity is Definition \ref{def:imagecore}(1) together with Definition \ref{def:sectioncore}(1). By Definition \ref{def:picosec}(2),
\[
\PicoSectionProj{\sigma}\circ\PicoMinSection{\sigma}
=
\mathrm{id}_{\NanoSectionBody{\sigma}}.
\]
Restricting \(\PicoSectionProj{\sigma}\) to \(\PicoImgBody{\sigma}\) yields
\[
\PicoImgProj{\sigma}\circ\PicoMinSection{\sigma}
=
\mathrm{id}_{\NanoSectionBody{\sigma}}.
\]
Conversely, every \(y\in\PicoImgBody{\sigma}\) has the form
\[
y=\PicoMinSection{\sigma}(p)
\qquad
\text{for some }
p\in\NanoSectionBody{\sigma}.
\]
Hence
\[
\PicoMinSection{\sigma}\bigl(\PicoImgProj{\sigma}(y)\bigr)
=
\PicoMinSection{\sigma}\bigl(\PicoSectionProj{\sigma}(y)\bigr)
=
\PicoMinSection{\sigma}\bigl(\PicoSectionProj{\sigma}(\PicoMinSection{\sigma}(p))\bigr)
=
\PicoMinSection{\sigma}(p)
=
y.
\]
Therefore \(\PicoImgProj{\sigma}\) and \(\PicoMinSection{\sigma}\) are mutually inverse continuous maps, so the restricted projection is a homeomorphism with inverse \(\PicoMinSection{\sigma}\).
\end{proof}

\begin{proposition}[The canonical minimizer section reconstructs the image-core functional]
\label{prop:reconstructfunctional}
For every benchmark-faithful pico-section package, the restricted image-core functional is recovered from the canonical minimizer section and the fixed nano-section functional by
\[
\PicoImgFunctional{\sigma}\bigl(\PicoMinSection{\sigma}(p)\bigr)
=
\NanoSectionFunctional{\sigma}(p)
\qquad
\text{for every }p\in\NanoSectionBody{\sigma}.
\]
Equivalently,
\[
\PicoImgFunctional{\sigma}(y)
=
\NanoSectionFunctional{\sigma}\bigl((\PicoMinSection{\sigma})^{-1}(y)\bigr)
\qquad
\text{for every }y\in\PicoImgBody{\sigma}.
\]
\end{proposition}

\begin{proof}
By Definition \ref{def:imagecore}(2),
\[
\PicoImgFunctional{\sigma}\bigl(\PicoMinSection{\sigma}(p)\bigr)
=
\PicoSectionFunctional{\sigma}\bigl(\PicoSectionNanoMin{\sigma}(p)\bigr).
\]
Definition \ref{def:picosec}(2) then gives
\[
\PicoSectionFunctional{\sigma}\bigl(\PicoSectionNanoMin{\sigma}(p)\bigr)
=
\NanoSectionFunctional{\sigma}(p)
\]
for every \(p\in\NanoSectionBody{\sigma}\). This proves the first displayed identity. The equivalent formula on \(\PicoImgBody{\sigma}\) follows from Proposition \ref{prop:reconstructbody}, which shows that every \(y\in\PicoImgBody{\sigma}\) has the unique form \(y=\PicoMinSection{\sigma}(p)\).
\end{proof}

\begin{proposition}[The shell image is exactly the shell restriction of the canonical minimizer section]
\label{prop:reconstructshell}
For every benchmark-faithful pico-section package,
\[
\PicoImgShell{\sigma}
=
\PicoMinSection{\sigma}\bigl(\NanoSectionShell{\sigma}\bigr),
\]
and the shell restriction
\[
\PicoMinSection{\sigma}\big|_{\NanoSectionShell{\sigma}}
:
\NanoSectionShell{\sigma}\xrightarrow{\cong}\PicoImgShell{\sigma}
\]
is a diffeomorphism with inverse
\[
\PicoImgProj{\sigma}\big|_{\PicoImgShell{\sigma}}
:
\PicoImgShell{\sigma}\xrightarrow{\cong}\NanoSectionShell{\sigma}.
\]
\end{proposition}

\begin{proof}
The displayed equality is Definition \ref{def:imagecore}(3) together with Definition \ref{def:sectioncore}(1). By Definition \ref{def:picosec}(3),
\[
\PicoSectionProj{\sigma}\big|_{\PicoSectionShell{\sigma}}
:
\PicoSectionShell{\sigma}\xrightarrow{\cong}\NanoSectionShell{\sigma}
\]
is a diffeomorphism, and
\[
\PicoMinSection{\sigma}\bigl(\NanoSectionShell{\sigma}\bigr)
\subset
\PicoSectionShell{\sigma}.
\]
Restricting the shell diffeomorphism to \(\PicoImgShell{\sigma}\) gives
\[
\PicoImgProj{\sigma}\big|_{\PicoImgShell{\sigma}}
\circ
\PicoMinSection{\sigma}\big|_{\NanoSectionShell{\sigma}}
=
\mathrm{id}_{\NanoSectionShell{\sigma}}.
\]
Conversely, if \(z\in\PicoImgShell{\sigma}\), then \(z=\PicoMinSection{\sigma}(p)\) for some \(p\in\NanoSectionShell{\sigma}\), so
\[
\PicoMinSection{\sigma}\bigl(\PicoImgProj{\sigma}(z)\bigr)
=
z.
\]
Thus the shell restriction of the canonical minimizer section and the restricted shell projection are inverse diffeomorphisms.
\end{proof}

\begin{proposition}[The full image core is reconstructed from the section core]
\label{prop:reconstructimage}
For every benchmark-faithful pico-section package, the full observer-relevant pico-section minimizer-image core \(\PicoImageCore\) is reconstructed from the smaller minimizer-section core \(\PicoMinSectionCore\) together with the fixed nano-section body \(\NanoSectionBody{\sigma}\), nano-section functional \(\NanoSectionFunctional{\sigma}\), and exact nano-section shell \(\NanoSectionShell{\sigma}\).
\end{proposition}

\begin{proof}
By Proposition \ref{prop:reconstructbody}, the image body is the image of the canonical minimizer section and the restricted projection is its inverse. By Proposition \ref{prop:reconstructfunctional}, the restricted image-core functional is uniquely determined by the canonical minimizer section together with the fixed nano-section functional. By Proposition \ref{prop:reconstructshell}, the shell image is exactly the shell restriction of that same section. Benchmark faithfulness is already part of Definition \ref{def:sectioncore}(2). Therefore each constituent of Definition \ref{def:imagecore} is recovered from Definition \ref{def:sectioncore} plus the fixed nano-section data.
\end{proof}

\begin{proposition}[All current image-core uses factor through the section core]
\label{prop:downstream}
For every benchmark-faithful pico-section package, all current benchmark-facing uses of the minimizer-image core \(\PicoImageCore\) on the active line factor through the minimizer-section core \(\PicoMinSectionCore\) together with the fixed nano-section data. In particular, the downstream recovery of the current micro-section package \(\MicroSectionPackage\), the current sub-section package \(\SubSectionPackage\), the observer-relevant sub-pre-core exact-shell transport-section core \(\SubPreSectionCore\), and the previously recorded shell-transport consequences does not require the separate image-core presentation once the section core is fixed.
\end{proposition}

\begin{proof}
Proposition \ref{prop:reconstructimage} reconstructs the full image core from the section core and the fixed nano-section body, functional, and shell. The active line uses the current image core only through those constituents together with the fixed nano-section downstream chain recorded in Definition \ref{def:downstream}(3). Therefore every current benchmark-facing use of \(\PicoImageCore\) factors through \(\PicoMinSectionCore\) together with the fixed nano-section data, and the separate image-core presentation is not needed to recover the recorded downstream consequences.
\end{proof}

\begin{proposition}[Separate image-core presentation is benchmark neutral at current scope]
\label{prop:neutral}
For every benchmark-faithful pico-section package, the separate presentation of
\[
\PicoImgBody{\sigma},
\qquad
\PicoImgProj{\sigma},
\qquad
\PicoImgFunctional{\sigma},
\qquad
\PicoImgShell{\sigma}
\]
does not add independent benchmark-facing freedom beyond the minimizer-section core \(\PicoMinSectionCore\).
\end{proposition}

\begin{proof}
By Proposition \ref{prop:reconstructimage}, each listed image-core constituent is reconstructed from the canonical minimizer section together with the fixed nano-section data. By Proposition \ref{prop:downstream}, all current benchmark-facing uses of the image core already factor through that reconstructed data. Therefore the separate image-core presentation is not additional benchmark-facing information beyond the smaller minimizer-section core.
\end{proof}

\begin{proposition}[The benchmark package remains fixed]
\label{prop:benchmark}
Every observer-relevant pico-section minimizer-section core preserves the same benchmark one-metric package \eqref{eq:fixedoutputs}. In particular, the operative metric remains exactly \(\CompletedMetric\), the ordinary matter route is unchanged, there is no second operative metric, and the low-energy matter law is not enlarged.
\end{proposition}

\begin{proof}
This is part of benchmark faithfulness in Definition \ref{def:sectioncore}(2). The present note only removes presentational redundancy inside the current pico-section image core. It does not alter the benchmark exact-closure package.
\end{proof}

\section{Main Theorem}

\begin{theorem}[Observer-relevant pico-section minimizer-image core collapse to an observer-relevant pico-section minimizer-section core]
\label{thm:main}
Fix the primitive continuation data of Definition \ref{def:fixed}, the recorded nano-section downstream chain of Definition \ref{def:downstream}, the recorded nano-section package of Definition \ref{def:nanosec}, and a benchmark-faithful pico-section package in the sense of Definition \ref{def:picosec}. Then, for each sector \(\sigma\in\{\Car,\Cur,\Hom\}\), the live benchmark-facing burden inside the current pico-section minimizer-image-core frontier collapses from the observer-relevant pico-section minimizer-image core \(\PicoImageCore\) to the smaller observer-relevant pico-section minimizer-section core \(\PicoMinSectionCore\). More precisely:
\begin{enumerate}
    \item the image body and restricted projection are reconstructed from the canonical minimizer section by Proposition \ref{prop:reconstructbody};
    \item the restricted image-core functional is reconstructed from the same section together with the fixed nano-section functional by Proposition \ref{prop:reconstructfunctional};
    \item the exact-shell image is exactly the shell restriction of that section by Proposition \ref{prop:reconstructshell};
    \item the full image core is therefore recovered from the section core and the fixed nano-section data by Proposition \ref{prop:reconstructimage};
    \item all current benchmark-facing uses of the image core factor through the section core by Proposition \ref{prop:downstream};
    \item the separate image-core presentation adds no independent benchmark-facing freedom beyond the section core by Proposition \ref{prop:neutral};
    \item the benchmark boundary remains fixed by Proposition \ref{prop:benchmark};
    \item consequently the remaining honest burden below the previous image-core frontier is no longer the full observer-relevant pico-section minimizer-image core \(\PicoImageCore\) as such. What remains is to derive or justify the sharper minimizer-section core \(\PicoMinSectionCore\) from still more primitive substrate structure, or else prove a still smaller theorem-level collapse strictly inside that section core.
\end{enumerate}
\end{theorem}

\begin{proof}
Item (1) is Proposition \ref{prop:reconstructbody}. Item (2) is Proposition \ref{prop:reconstructfunctional}. Item (3) is Proposition \ref{prop:reconstructshell}. Item (4) is Proposition \ref{prop:reconstructimage}. Item (5) is Proposition \ref{prop:downstream}. Item (6) is Proposition \ref{prop:neutral}. Item (7) is Proposition \ref{prop:benchmark}. Item (8) is the routing consequence of the previous items together with the active safeguard against counting another same-output deeper-origin relay as new front-line progress when the live benchmark-relevant datum is unchanged.
\end{proof}

\begin{corollary}[What must be done next]
\label{cor:next}
After Theorem \ref{thm:main}, the next honest benchmark-facing manuscript must either:
\begin{enumerate}
    \item derive or justify the observer-relevant pico-section minimizer-section core \(\PicoMinSectionCore\) from still more primitive substrate structure in a way that changes benchmark-facing status;
    \item identify a genuinely smaller theorem-level observer-relevant datum strictly inside \(\PicoMinSectionCore\) and prove a sharper collapse to it;
    \item do either of the previous two while preserving the same benchmark one-metric package and the same claim boundary.
\end{enumerate}
Another same-output deeper-origin relay that preserves this section core, another re-expansion back to the larger minimizer-image-core presentation as if its recovered constituents were still independently live, another reopening of the already-screened full pico-section package or the larger transport-image, sub-pre-core, or pre-core packages as if they were again the active burden, or any move that introduces a second operative metric, enlarges the ordinary matter law, or uses stronger promotion language does not satisfy the present burden. If a quantum continuation is pursued, it matters benchmark-facingly only insofar as it derives this same section core, collapses it further, or closes a benchmark derivation gate in a genuinely new way.
\end{corollary}

\begin{proof}
Theorem \ref{thm:main} shows that the image-core presentation is no longer the minimal benchmark-facing datum at this stage. Therefore a subsequent manuscript must improve on the smaller section core rather than merely restating the larger image core.
\end{proof}

\section{Conclusion}

The positive result of this note is a disciplined second compression step inside the current pico-section frontier. The previous active-primary theorem already showed that the full pico-section package is not the minimal benchmark-facing datum, because its benchmark-relevant content collapses to an observer-relevant minimizer-image core. The present theorem then takes the honest next internal route available there: it shows that even that image-core presentation is still larger than necessary.

What survives as the live burden is the observer-relevant pico-section minimizer-section core \(\PicoMinSectionCore\). Relative to the already fixed nano-section body, nano-section functional, and exact nano-section shell, that sharper datum is just the canonical minimizer section over the recorded nano-section body. The image body is its image, the exact-shell image is its shell restriction, the restricted projection is its inverse, and the restricted image-core functional is transported from the recorded nano-section functional along the same section. Those separately listed image-core constituents therefore do not add independent benchmark-facing freedom.

This is the honest next burden after the present theorem. A new primary manuscript must derive or justify that sharper section core from still more primitive substrate structure, or else collapse it further, while preserving the same benchmark one-metric package and the same claim boundary. The current result does not yet complete a first-principles derivation of GR from substrate microphysics, but it does narrow the live derivational burden one level further on the active line.

\end{document}
